\section{Current Results and Insights}

\subsection{Main OLMo-7B HFSA Grid}

The completed 3-seed grid supports a depth-scaling comparison between formal
logic and one-to-one natural-language traces. The current key result is not
that logic dominates everywhere, but that logic reaches stronger intermediate
depth extrapolation earlier. At the broadest train range, natural-language
exact traces narrow the gap or slightly catch up on strict joint validity.

\begin{center}
\begin{tabular}{lcc}
\toprule
Train range and eval point & Logic & NL exact \\
\midrule
Train 1--25, eval depth 50, joint@16 & 0.417 & 0.427 \\
\bottomrule
\end{tabular}
\end{center}

\subsection{Training Sequence Lengths}

\begin{center}
\begin{tabular}{lrrrr}
\toprule
Train range & Logic target & NL target & Logic total & NL total \\
\midrule
1--5 & 322 & 382 & 697 & 757 \\
1--10 & 500 & 653 & 1166 & 1319 \\
1--15 & 681 & 925 & 1637 & 1881 \\
1--20 & 863 & 1196 & 2109 & 2442 \\
1--25 & 1049 & 1469 & 2587 & 3008 \\
\bottomrule
\end{tabular}
\end{center}

\subsection{OOD Evaluation}

Bare-format OOD reruns are complete for the main OLMo grid. NL transfers much
better to GSM8K numeric exact match. Logic transfers much better to
context-provided HotpotQA, 2WikiMultiHopQA, and MuSiQue EM/F1. The current
interpretation is that the QA tasks measure context-QA robustness under the
learned answer format; they do not yet establish explicit reasoning-trace
quality.

\subsection{Scratch Pretraining}

Tiny Llama 20k and 100k scratch-pretraining runs are complete. They show useful
train-band behavior and serve as mechanism smoke tests, but strict depth-50
joint validity remains essentially absent.

\subsection{Architecture Ablations}

Qwen-2.5-7B, Qwen-2.5-1.5B, Gemma-3-4B, and an OLMo-2-32B short-context pilot
are complete and tabulated in the generated analysis artifacts. These should be
folded into the manuscript tables and architecture comparison figures once the
current ablation wave finishes.
